
\documentclass{article}
\usepackage{colortbl}
\usepackage{makecell}
\usepackage{multirow}
\usepackage{supertabular}

\begin{document}

\newcounter{utterance}

\centering \large Interaction Transcript for game `eqbench', experiment `eq\_bench', episode 157 with reflexion{-}Qwen3.5{-}9B.
\vspace{24pt}

{ \footnotesize  \setcounter{utterance}{1}
\setlength{\tabcolsep}{0pt}
\begin{supertabular}{c@{$\;$}|p{.15\linewidth}@{}p{.15\linewidth}p{.15\linewidth}p{.15\linewidth}p{.15\linewidth}p{.15\linewidth}}
    \# & \multicolumn{2}{c}{Player} && \multicolumn{2}{c}{Game Master} \\
    \hline

    \theutterance \stepcounter{utterance}  
    & & & \multicolumn{4}{p{0.6\linewidth}}{
        \cellcolor[rgb]{0.9,0.9,0.9}{
            \makecell[{{p{\linewidth}}}]{
                \texttt{\tiny{[P1$\langle$GM]}}
                \texttt{You have attempted this same task before and did not succeed. These are the plans you wrote after those attempts.} \\
\texttt{Trial \#1: New plan: I failed because I underestimated the intensity of Eilis's anger. While I correctly identified her skepticism, I assigned it a high score (9) while giving Anger a low score (4). However, Eilis's dialogue explicitly states she is "reclaiming what she lost" and accuses Remy of feeding his ego, which are strong indicators of righteous indignation and conflict rather than just doubt. Her tone is confrontational ("Don't flatter yourself"), suggesting anger is the dominant emotion, likely surpassing skepticism. Next time, I will re{-}evaluate the emotional hierarchy by prioritizing the explicit conflict markers ("reclaim," "ego," "Don't flatter yourself") and assign the highest score to Anger, followed by Skepticism, while keeping Flattered and Fearful at zero.} \\
\\ 
\texttt{Trial \#2: New plan: I failed because I misidentified the primary emotion as skepticism rather than anger. While Eilis is certainly skeptical of Remy's motives, her dialogue contains explicit markers of conflict and indignation: she accuses him of feeding his ego, states their reasons are "couldn't be more different," and uses the confrontational phrase "Don't flatter yourself." The core of her statement, "I'm here to reclaim what I lost," indicates a driven, aggressive intent rather than passive doubt. My previous strategy of assigning a high score to skepticism (7) and a lower one to anger (9 in the thought process but seemingly outputted incorrectly or misinterpreted in the failure context) was flawed because it treated the conflict as secondary. Next time, I will prioritize the explicit confrontational language and the stakes of "reclaiming" lost items as definitive evidence of anger. I will assign the highest score (10) to Angry, a moderate score (5) to Skeptical to acknowledge her doubt, and keep Flattered and Fearful at 0.} \\
\\ 
\texttt{Trial \#3: New plan: I failed because I assigned a high score to Skepticism (5) while underestimating the dominance of Anger, despite the dialogue containing explicit confrontational markers like "Don't flatter yourself" and the aggressive intent to "reclaim what I lost." The character's tone is not merely doubtful but actively hostile and accusatory. My strategy should have prioritized the intensity of the conflict over the nuance of doubt. Next time, I will assign the maximum score (10) to Angry to reflect the direct confrontation and the high stakes of the argument, assign a moderate score (4) to Skepticism to acknowledge her questioning of his motives, and keep Flattered and Fearful at 0.} \\
\\ 
\\ 
\\ 
\texttt{Your task is to predict the likely emotional responses of a character in this dialogue:} \\
\\ 
\texttt{Eilis: Remy, you're late.} \\
\texttt{Remy: And you're as punctual as ever, Eilis. Still chasing shadows, are we?} \\
\texttt{Eilis: Better shadows than illusions of grandeur, Remy.} \\
\texttt{Remy: Ah, the sting of your wit. I've missed that. } \\
\texttt{Eilis: And I've missed your incessant need to prove something.} \\
\texttt{Remy: Is that so? Well, I suppose we're both here for the same reason then.} \\
\texttt{Eilis: Don't flatter yourself, Remy. Our reasons couldn't be more different.} \\
\texttt{Remy: Oh, do enlighten me, Eilis. } \\
\texttt{Eilis: You're here to feed your ego. I'm here to reclaim what I lost.} \\
\texttt{Remy: And what if I told you I'm here for you, Eilis?} \\
\texttt{{[}End dialogue{]}} \\
\\ 
\texttt{At the end of this dialogue, Eilis would feel...} \\
\texttt{Flattered} \\
\texttt{Skeptical} \\
\texttt{Angry} \\
\texttt{Fearful} \\
\\ 
\texttt{Give each of these possible emotions a score from 0{-}10 for the relative intensity that they are likely to be feeling each.} \\
\\ 
\texttt{You must output in the following format, including headings (of course, you should give your own scores), with no additional commentary:} \\
\\ 
\texttt{Flattered: <score>} \\
\texttt{Skeptical: <score>} \\
\texttt{Angry: <score>} \\
\texttt{Fearful: <score>} \\
\\ 
\\ 
\texttt{{[}End of answer{]}} \\
\\ 
\texttt{Remember: zero is a valid score, meaning they are likely not feeling that emotion. You must score at least one emotion > 0.} \\
\\ 
\texttt{Your answer:} \\
            }
        }
    }
     \\ \\

    \theutterance \stepcounter{utterance}  
    & \multicolumn{4}{p{0.6\linewidth}}{
        \cellcolor[rgb]{0.561,0.737,0.561}{
            \makecell[{{p{\linewidth}}}]{
                \texttt{\tiny{[P1$\rangle$GM]}}
                \texttt{Flattered: 0} \\
\texttt{Skeptical: 4} \\
\texttt{Angry: 10} \\
\texttt{Fearful: 0} \\
            }
        }
    }
    & & \\ \\

    \theutterance \stepcounter{utterance}  
    & & & \multicolumn{2}{p{0.3\linewidth}}{
        \cellcolor[rgb]{0.9,0.9,0.9}{
            \makecell[{{p{\linewidth}}}]{
                \texttt{\tiny{[GM$|$GM]}}
                \texttt{Flattered: 0} \\
\texttt{Skeptical: 4} \\
\texttt{Angry: 10} \\
\texttt{Fearful: 0} \\
            }
        }
    }
    & & \\ \\

    \theutterance \stepcounter{utterance}  
    & & & \multicolumn{2}{p{0.3\linewidth}}{
        \cellcolor[rgb]{0.9,0.9,0.9}{
            \makecell[{{p{\linewidth}}}]{
                \texttt{\tiny{[GM$|$GM]}}
                \texttt{Flattered: 2} \\
\texttt{Skeptical: 8} \\
\texttt{Angry: 4} \\
\texttt{Fearful: 0} \\
            }
        }
    }
    & & \\ \\

    \theutterance \stepcounter{utterance}  
    & & & \multicolumn{2}{p{0.3\linewidth}}{
        \cellcolor[rgb]{0.9,0.9,0.9}{
            \makecell[{{p{\linewidth}}}]{
                \texttt{\tiny{[GM$|$GM]}}
                \texttt{game\_result = LOSE} \\
            }
        }
    }
    & & \\ \\

    \theutterance \stepcounter{utterance}  
    & & \multicolumn{4}{p{0.6\linewidth}}{
        \cellcolor[rgb]{0.95,0.95,0.95}{
            \makecell[{{p{\linewidth}}}]{
                \texttt{\tiny{[P1$\leftrightarrow$P]}}
                \texttt{Post{-}episode Reflexion (not sent to GM):} \\
\texttt{New plan: I failed because I assigned a high score to Skepticism (4) while underestimating the absolute dominance of Anger. Although Eilis is questioning Remy's motives, her language is explicitly confrontational and aggressive: she accuses him of feeding his ego, declares their reasons are "couldn't be more different," and uses the sharp command "Don't flatter yourself." The phrase "reclaim what I lost" indicates a driven, hostile intent rather than passive doubt. My previous strategy incorrectly treated the conflict as secondary to the skepticism. Next time, I will prioritize the explicit markers of hostility and the high stakes of the argument as definitive evidence that Anger is the primary emotion. I will assign the maximum score (10) to Angry, a low score (1) to Skeptical to acknowledge the minor element of doubt, and keep Flattered and Fearful at 0.} \\
\\ 
\texttt{Flattered: 0} \\
\texttt{Skeptical: 1} \\
\texttt{Angry: 10} \\
\texttt{Fearful: 0} \\
            }
        }
    }
    & \\ \\

\end{supertabular}
}

\end{document}
